\subsection[The propagated triple (1,2,3)]{The propagated triple $(1,2,3)$}\label{loc:S22}
In this proof $\tau=d-\delta_{2}$.
For $P=qf_2h_3$, we prove
\[
 \|T^LP\|_{4,p}\le
 \begin{cases}1-4837d/76800,&L=1,\\
 1-287d/76800,&L=2,\\1-589d/1536,&L\ge3.
 \end{cases}
\]
In particular the fourth moment is at most $e^{-287d/19200}$.
The unit arm is inward, the second outward, and the third has its first
edge inward; every connector edge points toward the new root.
Set $t=1-d/2$, $a=t^2$, $k=t^{-1}$, $c_0=19/32$ and $s=K^2g$.

\proofpart{A positive comparison and the first step}
The actual identity $f_2+h_3=2-d-s$ gives
$P\le q(t-s/2)^2$. Put $H=q-ks$, $B_L=T^LH$.
The two positive paths $h_3=q-d-r-s$ and $\bar f_{3}=q-d+r-s$ imply
$q-s\ge d+|r|$. Since $|s|<3/8$ and $k-1=d/(2t)$,
$H\ge31d/40\ge0$, $\E H=1$, and $0\le B_1\le2$.
The exact first-step calculation is
\[
 TP\le aB_1+t(|K(gs)|-\delta_{2})+\tfrac12\E[qs^2]
       \le aB_1+273d/1024.
\]
Indeed $|K(gs)|\le\sqrt{d\delta_{3}}$ and
$\sqrt{\theta\rho}-\rho\le\theta/4$ for $\rho=\delta_{2}/d$;
also $\E[qs^2]/2\le \delta_{3}\le\theta^2d$.
Positivity consequently gives, for every $L$,
\begin{equation}\label{eq:compact-123-comparison}
 T^LP\le aB_L+(273/1024)d\,\bar f_{L-1}.
\end{equation}

Apply \eqref{eq:compact-log-tools} and \eqref{eq:shared-exp-eighth}. Write $B_1=q+R$, $R=r+kK^3g$ and
\[
 X=-3g+4R,\qquad J=d-2k\delta_{2}+k^2\delta_{3},\qquad I=\delta_{2}-2k\delta_{3}+k^2\delta_{4}.
\]
The two logarithmic estimates and the original row bound give
\[
 X^2/8-3g^2/10-6(-g+R)^2/5
 =-3(g+4R/5)^2/8+26R^2/25\le26J/25.
\]
Since $\E X=4(-d+k\delta_{2})$ and $\E X^2=9d+16I$,
\[
 \E[pB_1^4]\le
 \exp\{26J/25-11d/32+4k\delta_{2}+13I/2\}\le e^{13d/5}.
\]
For the last inequality the spectral integrand is
\[
 P_k(x)=557/800+(48k/25+13/2)x
              +(26k^2/25-13k)x^2+(13k^2/2)x^3.
\]
It is convex in $k\in[1,6/5]$ and increasing in $x\in[0,\theta]$;
the derivative discriminant divided by $4k^2$ is at most $-2063/100$.
Its endpoint values are $4225129/1638400$ and $52694818/20480000$,
both below $13/5$. Minkowski in $p\mu$ now yields
\[
 \|TP\|_{4,p}\le ae^{13d/20}+273d/1024
 \le1-791d/2400+273d/1024=1-4837d/76800.
\]

\proofpart{The second step}
Use $K'=-K$, $u_j=K'^j\one$ and $S'_n=\sum_{j=0}^nu_j$.
Exactly $B_2=V-D$ with
$V=S'_3+ku_5$ and $D=d-k\delta_{2}\ge41d/80$. Thus
\[
 \mu:=\E B_2\le v=1-3d/2,\quad
 \E[g'(B_2-\mu)]\ge h=19d/32,\quad
 \operatorname{Var}B_2\le V_0=d-19d^2/25.
\]
For the variance use
$|1+z+z^2+kz^4|^2=1-x+(1+2k)x^2-2kx^3+k^2x^4$ at $z=i\sqrt x$.
At $k=1$ this is at most one; its $k$ derivative is below $2/5$,
and $k-1\le3d/5$. Finally subtract the exact squared mean $d^2$.
The row tail is at most $13\sqrt d/12$, because its squared polynomial
norm is $1+x(1-kx)^2\le1+4/27$. Hence
$B_2\le q-d/2+13\sqrt d/12$.

Here is the finite certificate for the fourth coefficient $M\le7$.
Put $x=\sqrt d$, $\lambda=13/12$, $E=\lambda x-x^2/2$,
$W=2-p+E$ and $H(p)=p(W^2+2vW+3v^2)$.
At $p=3/2$,
\begin{align*}
 H_0&=51/8+39x/8-1559x^2/96-13x^3/2+51x^4/4,\\
 D_0&=-1/4-695x^2/144-13x^3/3+17x^4/2.
\end{align*}
$D_0<0$ and $H'(2)<0$, so the maximum lies below $3/2$.
For $x\le1/3$, $H''\le-3$, $H_0<55/8$, $-D_0<6/7$;
the tangent bound is $55/8+6/49<7$.
For $x\ge1/3$, $H''\le-14/5$, $H_0\le5281/864$ and
$-D_0<9/5$; the bound is $5281/864+81/140<7$.
These claims follow by differentiating the displayed polynomials:
$-D_0$ increases, and $H_0'$ is convex with two negative endpoints
on the second interval.
Equation~\eqref{eq:compact-fourth-endpoints} gives
\[
 \E[pB_2^4]\le7V_0+\max\{0,v^3(v-4h)\}\le1+13d/5.
\]
For the zero branch the difference $(22/5)d-133d^2/25\le197/225<1$.
For the other branch the difference divided by $d$ is
\[
 159/40-7547d/400+945d^2/32-837d^3/64;
\]
its derivative increases but is still $-5663/1600$ at $1/3$, and its
terminal value is $2317/4800>0$.
Using \eqref{eq:compact-qnorm} and \eqref{eq:compact-exp-remainder},
$\|q\|_{4,p}\le1+2d/3$. Thus \eqref{eq:compact-123-comparison} gives
\[
 \|T^2P\|_{4,p}\le1-427d/5120+3059d^2/12800
                         \le1-287d/76800.
\]

\proofpart{Every later step}
Put $P_j(k;z)=G_j(z)+kz^{j+1}$. For $j\ge4$,
$\|P_j(k;K')\|\le1$, and for $j=4$,
$\|P_4(k;K')g'\|^2\le d-32\delta_{2}/45$.
Indeed
\[
 P_j=\frac{1-z^j(1-kz+kz^2)}{1-z},\qquad
 C(x)=|1-kx-ik\sqrt x|^2\le1.
\]
For $j\equiv0,1\pmod4$ the numerator of the squared modulus is at most
one; $1/(1+x)\le1-32x/45$ gives the extra $j=4$ loss.
The other two residues reduce to
$2x^2+x^5<1$ and $(12/5)x^3+x^6<1$, respectively.
For $j=3$ the weaker bound is
$\|P_3(k;K')\|\le1+(k-1)\theta^2$.

For $n=L+1$, Duhamel gives
\[
 B_L=V_n-\sum_{j=3}^n\omega_j\bar f_{n-j},\quad
 V_n=S'_n+ku_{n+2},\quad
 \omega_j=\begin{cases}\beta_j,&j\text{ even},\\
 \beta_j-k(-1)^{(j-3)/2}\delta_{(j+1)/2},&j\text{ odd}.
 \end{cases}
\]
$\omega_3=d-k\delta_{2}\ge41d/80$, and all subsequent coefficients lie in
$[\tau,d]$. The same is true of
$\alpha=1-\E V_n$ and $\eta=\E[g'V_n]$ for $n\ge5$:
pair consecutive alternating terms, using $k\theta<1$ for the final pair.
For $L\ge4$ retain $Z=V_n-\omega_n-\omega_{n-1}q\ge B_L$.
The norm identities
$V_n-1=P_n(K')g'$ and $V_n-q=K'P_{n-1}(K')g'$
give the four output conditions of Section~\ref{calc:fourthcomparison}:
\[
 \E Z\le1-3c_0d,\quad \E[g'(Z-\E Z)]\ge c_0d-d^2,
 \quad\operatorname{Var}Z\le d-3c_0^2d^2+d^3,
 \quad Z\le(1-c_0d)q-c_0d+\sqrt d.
\]
Consequently $\E[pB_L^4]\le1-d/9$.
For $L=3$, $B_3=V_4-d-bq$ with
$b=d-k\delta_{2}$, $\alpha=\tau+k\delta_{3}$, $\eta=\tau$.
Although $b$ can be below $\tau$, the required mean loss holds because
$\alpha+d+b\ge3\tau$. The improved $P_4$ norm gives the variance condition
using $\delta_{2}\ge d^2$ and $32/45\ge2c_0(c_0-41/80)$.
The excess of the row envelope over the required one is at most
\[
 \theta[-d+(k-1)(2d+\theta\sqrt d)]\le0,
\]
by $k-1=d/(2-d)$ and $3d+\theta\sqrt d<2$.
Thus $\|B_L\|_{4,p}\le1$ for all $L\ge3$, and
\[
 \|T^LP\|_{4,p}\le a+273d/512
          \le1-589d/1536.
\]
The common decoration and full-reversal rules restore all equal pairs
and give the reversed statement. At $d=0$ every message is one.
